\chapter{Momentum Sudut Kuantum}\label{ch:b3:quantum-angular-momentum}

Arahkanlah teleskop radio ke sepetak langit yang gelap lalu talalah pada
\qty{115}{GHz}: muncul sebuah garis terang setajam jarum --- yakni
molekul karbon monoksida, sepuluh kelvin di atas nol mutlak, yang
menuruni satu anak tangga dari tangga keadaan \emph{rotasi}. Jungkir
baliknya, seperti segala sesuatu yang berputar dalam mekanika kuantum,
terkuantumkan dua kali: besar momentum sudutnya hanya boleh
$\sqrt{l(l+1)}\,\hbar$, dan komponennya sepanjang sumbu mana pun yang
dipilih hanya boleh $m\hbar$. Bab ini menurunkan kuantisasi ganda itu ---
sekali secara aljabar, dari \omterm{def:b3:quantum-formalism:commutator}{komutator} yang diwarisi dari \omterm{def:b3:hamiltonian-mechanics:poisson}{kurung Poisson}
pada \cref{ch:b3:hamiltonian-mechanics}, dan sekali secara nyata,
sebagai \omterm{prop:b3:quantum-angular-momentum:harmonics}{harmonik bola} yang membentuk tiap orbital atom. Aljabarnya akan
menyerahkan lebih daripada yang kita minta: ia mengizinkan nilai
setengah bulat yang tak dapat diwujudkan gerak orbital mana pun, sebuah
kekosongan yang diisi alam dua bab dari sini dengan spin. Di
perjalanannya kita bertemu tangga rotasi molekul --- garis gelombang
milimeter yang dipakai para astronom untuk menimbang gas dingin galaksi
--- dan momen magnetik yang lewatnya momentum sudut pertama kali
menampakkan diri dalam keadaan terbelah.

\section{Aljabar rotasi}

\begin{definition}[Operator momentum sudut]\label{def:b3:quantum-angular-momentum:operators}
\index{momentum sudut!kuantum}
Momentum sudut orbital adalah operator $\hat{\vect L} =
\hat{\vect r}\wedge\hat{\vect p}$, dengan komponen $\hat L_z = \hat
x\hat p_y - \hat y\hat p_x$ beserta pasangan sikliknya. Dari $[\hat x, \hat p_x] =
\iu\hbar$:
\[
[\hat L_x, \hat L_y] = \iu\hbar\,\hat L_z
\quad\text{(dan siklik)} , \qquad
[\hat L^2, \hat L_z] = 0 :
\]
komponennya saling tak serasi --- tak ada keadaan yang mempunyai dua di
antaranya tajam --- tetapi kuadrat totalnya serasi dengan salah satu di
antaranya: pasangan $(\hat L^2, \hat L_z)$ adalah pilihan baku bagi
labelnya. \omterm{def:b3:quantum-formalism:commutator}{Komutator} ini tepat $\iu\hbar$ kali \omterm{def:b3:hamiltonian-mechanics:poisson}{kurung Poisson} yang
dihitung dalam \cref{ex:b3:hamiltonian-mechanics:brackets}: yakni kamus
Dirac yang sedang bekerja.
\end{definition}

\begin{theorem}[Spektrumnya, hanya dari aljabarnya]\label{thm:b3:quantum-angular-momentum:spectrum}
\index{bilangan kuantum magnetik}
Misalkan $\hat J_x, \hat J_y, \hat J_z$ adalah tiga \omterm{def:b3:quantum-formalism:observable}{operator Hermitian}
\emph{sembarang} yang memenuhi hubungan komutasi di atas. Maka \omterm{def:b3:quantum-formalism:observable}{nilai
eigen} bersama $(\hat J^2, \hat J_z)$ adalah
\[
\hat J^2:\ j(j+1)\hbar^2 , \qquad
\hat J_z:\ m\hbar , \quad m = -j, -j+1, \dots, +j ,
\]
dengan $j$ adalah \emph{bilangan bulat atau setengah bulat} yang tak
negatif: ada $2j + 1$ nilai $m$ untuk tiap $j$. \omterm{def:b3:harmonic-oscillator:ladder}{Operator tangga}
$\hat J_\pm = \hat
J_x \pm \iu\hat J_y$ melangkahkan $m$ sebesar $\pm1$ pada $j$ tetap:
\[
\hat J_\pm\ket{j,m} = \hbar\sqrt{j(j+1) - m(m\pm1)}\,\ket{j,m\pm1} .
\]
\end{theorem}

\begin{proof}[Bukti sebagian]
$[\hat J_z, \hat J_\pm] = \pm\hbar\hat J_\pm$: mengenakan $\hat
J_\pm$ menggeser \omterm{def:b3:quantum-formalism:observable}{nilai eigen} $\hat J_z$ sebesar $\pm\hbar$, pada
$\hat J^2$ tetap (yang berkomutasi dengan segala yang dibangun dari $\hat
J_i$). Normanya $\|\hat J_\pm\ket{j,m}\|^2 = \hbar^2[j(j+1) -
m(m\pm1)]$ (yang dihitung dari $\hat J_\mp\hat J_\pm = \hat J^2 - \hat
J_z^2 \mp \hbar\hat J_z$) harus tetap tak negatif: tangganya karena itu
harus berhenti di atas pada suatu $m_{\max}$ dengan $m_{\max}(m_{\max}+1) =
j(j+1)$, yakni $m_{\max} = j$, dan di bawah pada $m_{\min} = -j$.
Memanjat dari $-j$ ke $+j$ dalam langkah satuan memaksa $2j$ menjadi
bilangan bulat tak negatif. Penulisan \omterm{def:b3:quantum-formalism:observable}{nilai eigennya} $j(j+1)\hbar^2$
dengan demikian dibenarkan setelah kejadiannya.
\end{proof}

\begin{remark}[Sebuah kekosongan dalam katalognya]\label{rem:b3:quantum-angular-momentum:halfinteger}
Aljabarnya mengizinkan $j = \tfrac12, \tfrac32, \dots$ --- yakni tangga
dengan anak tangga berjumlah genap. Gerak orbital, seperti akan kita
tunjukkan berikut ini, hanya memakai bilangan bulat. Namun alam tak
membuang apa pun: elektron itu sendiri membawa $j = \tfrac12$, tanpa
orbit di belakangnya (\cref{ch:b3:spin-two-level}); aljabar yang
diturunkan di sini, tanpa perubahan, akan menjalankan kemagnetan atom,
spin inti dan kubit.
\end{remark}

\section{Momentum sudut orbital: harmonik bola}

\begin{proposition}[Harmonik bola]\label{prop:b3:quantum-angular-momentum:harmonics}
\index{harmonik bola}\index{bilangan kuantum orbital}
Dalam koordinat bola $\hat L_z = -\iu\hbar\,\partial_\varphi$,
dan fungsi eigen bersama $(\hat L^2, \hat L_z)$ pada bolanya adalah
\emph{\omterm{prop:b3:quantum-angular-momentum:harmonics}{harmonik bola}} $Y_l^m(\theta, \varphi)$:
\[
\hat L^2\,Y_l^m = l(l+1)\hbar^2\,Y_l^m , \qquad
\hat L_z\,Y_l^m = m\hbar\,Y_l^m ,
\]
dengan $l = 0, 1, 2, \dots$ \emph{bilangan bulat} --- ketunggalan nilai
$\eu^{\iu m\varphi}$ memaksa $m$ bulat, sehingga $l$ pun bulat. Beberapa
yang pertama, kecuali penormalannya: $Y_0^0 = \text{const}$ (bentuk $s$,
sebuah bola); $Y_1^0 \propto \cos\theta$ dan $Y_1^{\pm1}
\propto \sin\theta\,\eu^{\pm\iu\varphi}$ (bentuk $p$, dua cuping);
$Y_2^0 \propto 3\cos^2\theta - 1$ beserta pasangannya (keluarga $d$).
Paritasnya: $Y_l^m(-\vect r$ arah$) = (-1)^l\,Y_l^m$.
\end{proposition}
\admitted

\begin{omfigure}
\begin{tikzpicture}[font=\small, scale=0.94]
  % s orbital
  \begin{scope}
    \draw[->] (0,-1.7) -- (0,1.8);
    \node[anchor=south] at (0,1.8) {$z$};
    \draw[omDef, thick, fill=omDef!18] (0,0) circle (1.1);
    \node at (0,-2.3) {$l = 0$ ($s$): bola};
  \end{scope}
  % p_z
  \begin{scope}[xshift=4.4cm]
    \draw[->] (0,-1.7) -- (0,1.8);
    \node[anchor=south] at (0,1.8) {$z$};
    \draw[omThm, thick, fill=omThm!18, domain=0:360, smooth, samples=100]
      plot ({1.45*abs(cos(\x))*sin(\x)}, {1.45*abs(cos(\x))*cos(\x)});
    \node at (0,-2.3) {$l = 1$, $m = 0$: $|\cos\theta|$};
  \end{scope}
  % p_+-1
  \begin{scope}[xshift=8.8cm]
    \draw[->] (0,-1.7) -- (0,1.8);
    \node[anchor=south] at (0,1.8) {$z$};
    \draw[omExo, thick, fill=omExo!18, domain=0:360, smooth, samples=100]
      plot ({1.45*abs(sin(\x))*sin(\x)}, {1.45*abs(sin(\x))*cos(\x)});
    \node at (0,-2.3) {$l = 1$, $m = \pm1$: $|\sin\theta|$};
  \end{scope}
  % d_z2
  \begin{scope}[xshift=13.2cm]
    \draw[->] (0,-1.7) -- (0,1.8);
    \node[anchor=south] at (0,1.8) {$z$};
    \draw[omProp, thick, fill=omProp!18, domain=0:360, smooth, samples=140]
      plot ({0.62*abs(3*cos(\x)*cos(\x)-1)*sin(\x)}, {0.62*abs(3*cos(\x)*cos(\x)-1)*cos(\x)});
    \node at (0,-2.3) {$l = 2$, $m = 0$};
  \end{scope}
\end{tikzpicture}

{\small Diagram polar $|Y_l^m(\theta)|$ (penampang pada bidang yang
memuat sumbu $z$; bentuk lengkapnya adalah benda putarnya).
Kerangka sudut ini, yang tak bergantung pada potensial apa pun, akan
mendandani tiap atom dalam \cref{ch:b3:hydrogen-atom}.}
\end{omfigure}

\begin{example}[Rotor tegar dan tangganya]\label{ex:b3:quantum-angular-momentum:rotor}
Molekul dwiatom yang berjungkir balik dengan momen kelembaman $I$ mempunyai $\hat H
= \hat L^2/2I$: energinya
\[
E_J = \frac{\hbar^2}{2I}\,J(J+1) = B\,J(J+1) ,
\qquad J = 0, 1, 2, \dots
\]
masing-masing berdegenerasi $(2J+1)$ kali. Serapan fotonnya memenuhi $\Delta J =
\pm1$, sehingga frekuensi serapannya $\nu_{J+1\leftarrow J} =
2B(J+1)/h$: sebuah sisir garis yang \emph{berjarak sama} --- yakni ciri
yang membuat spektrum rotasi dikenali dengan sekali pandang. Untuk
karbon monoksida, $B/h = \qty{57.6}{GHz}$: garis dasarnya jatuh pada
\qty{115}{GHz} ($\lambda = \qty{2.6}{mm}$), yakni garis andalan astronomi
radio milimeter
(\cref{pb:b3:quantum-angular-momentum:1}).
\end{example}

\section{Potensial sentral, dilengkapi}

\begin{theorem}[Pemisahan dalam potensial sentral apa pun]\label{thm:b3:quantum-angular-momentum:radial}
\index{persamaan radial}\index{penghalang sentrifugal}
Untuk $V(r)$, keadaan stasionernya dapat diambil sebagai
\[
\psi(\vect r) = \frac{u(r)}{r}\,Y_l^m(\theta, \varphi) ,
\]
dengan $u$ menyelesaikan \omterm{prop:b3:schrodinger-three-dimensions:swave}{persamaan radial} satu dimensi
\[
-\frac{\hbar^2}{2m}\,u'' + \Big[V(r) +
\frac{\hbar^2\,l(l+1)}{2mr^2}\Big]u = E\,u , \qquad u(0) = 0 .
\]
Soal sudutnya diselesaikan sekali untuk selamanya oleh $Y_l^m$; tiap
$l$ menambahkan \emph{\omterm{thm:b3:quantum-angular-momentum:radial}{penghalang sentrifugal}} yang menolak
$\hbar^2l(l+1)/2mr^2$ --- yakni versi kuantum potensial efektif pada
\cref{ex:b3:lagrangian-mechanics:central} --- dan tiap aras dengan $l$
tertentu berdegenerasi $(2l+1)$ kali, karena tak ada gaya sentral yang
peduli pada arah sumbu $z$. Muslihat gelombang $s$ pada
\cref{prop:b3:schrodinger-three-dimensions:swave} adalah baris
$l = 0$ teorema ini.
\end{theorem}

\begin{proof}[Bukti sebagian]
Laplaciannya terbelah sebagai $\Delta = \tfrac1r\partial_r^2(r\,\cdot) -
\hat L^2/\hbar^2r^2$ (diterima tanpa bukti; itulah pernyataan bahwa
$\hat L^2$ merupakan bagian sudut $-\hbar^2\Delta$). Masukkanlah
$\psi = (u/r)Y_l^m$ lalu pakailah \omterm{def:b3:quantum-formalism:observable}{nilai eigen} $\hat L^2$: diperolehlah
persamaan yang dinyatakan itu. Degenerasinya menyusul karena $\hat H$
berkomutasi dengan ketiga $\hat
L_i$, yang \omterm{def:b3:harmonic-oscillator:ladder}{operator tangganya} menggeser
$m$ tanpa mengubah energinya.
\end{proof}

\begin{omfigure}
\begin{tikzpicture}
  \begin{axis}[omaxis, axis x line=bottom, axis y line=left,
    width=10.6cm, height=6.0cm, xmin=0, xmax=8, ymin=-1.35, ymax=1.3,
    xlabel={$r$ (dalam satuan $a_0$)}, ylabel={$U_{\text{eff}}$ (dalam satuan $E_{\text{I}}$)},
    xlabel style={at={(axis description cs:0.5,-0.1)}, anchor=north},
    xtick={2,4,6}, ytick=\empty,
    legend style={font=\scriptsize, at={(0.97,0.95)}, anchor=north east,
      draw=none, fill=none}]
    \addplot[omDef, thick, domain=1.35:8, samples=120] {-2/x};
    \addlegendentry{$l = 0$: Coulomb murni}
    \addplot[omThm, thick, domain=0.42:8, samples=160] {-2/x + 1/(x*x)};
    \addlegendentry{$l = 1$}
    \addplot[omExo, thick, domain=0.85:8, samples=160] {-2/x + 3/(x*x)};
    \addlegendentry{$l = 2$}
    \draw[gray, dashed] (axis cs:0,0) -- (axis cs:8,0);
  \end{axis}
\end{tikzpicture}

{\small Potensial radial efektif soal Coulomb: \omterm{thm:b3:quantum-angular-momentum:radial}{penghalang sentrifugal}
$\hbar^2l(l+1)/2mr^2$ menembok titik asalnya untuk
$l \ge 1$. Hanya keadaan $s$ yang menyentuh intinya --- dengan akibat
mulai dari spektrum atom sampai penangkapan elektron radioaktif.}
\end{omfigure}

\section{Momen magnetik: momentum sudut yang menjadi kasat mata}

\begin{proposition}[Momen magnetik orbital dan efek Zeeman]\label{prop:b3:quantum-angular-momentum:magnetic}
\index{magneton Bohr}\index{efek Zeeman}
Sebuah partikel bermuatan $q$ dan bermassa $m$ dengan momentum sudut
orbital $\hat{\vect L}$ membawa momen magnetik
\[
\hat{\vect\mu} = \frac{q}{2m}\,\hat{\vect L} ;
\]
untuk elektron, satuan wajarnya adalah \emph{\omterm{prop:b3:quantum-angular-momentum:magnetic}{magneton Bohr}}
$\mu_{\text{B}} = e\hbar/2m_{\text{e}} = \qty{5.79e-5}{eV/T}$
(bandingkanlah \cref{exo:b3:schrodinger-three-dimensions:5}). Dalam
medan $B\vect e_z$ energi $-\hat{\vect\mu}\cdot\vect B$ menggeser tiap
arasnya sebesar $+m\,\mu_{\text{B}}B$ (muatan elektronnya negatif):
sebuah aras dengan $l$ tertentu terbelah menjadi $2l + 1$ komponennya
--- yakni \emph{\omterm{prop:b3:quantum-angular-momentum:magnetic}{efek Zeeman}}, peragaan langsung pertama kuantisasi
$m$, sekaligus alasan $m$ disebut \omterm{thm:b3:quantum-angular-momentum:spectrum}{bilangan kuantum magnetik}.
Pada $B = \qty{1}{T}$ pembelahannya $\qty{58}{\micro eV}$: kecil
dibandingkan energi optis, namun mudah dipilah sebagai pergeseran garis
spektrum --- dan, dibaca terbalik, ia sebuah magnetometer: pembelahan
Zeeman garis sinar matahari memetakan medan magnet bintik matahari.
\end{proposition}

\begin{proof}[Bukti sebagian]
Secara klasik, orbit sebuah muatan adalah lingkar arus:
$\mu = IA = (qv/2\pi r)(\pi r^2) = qL/2m$; pernyataan
operatornya mewarisinya
(dan menyusul dari kopling $\hat{\vect A}$ pada
\cref{prop:b3:lagrangian-mechanics:charge}). Pergeseran energinya adalah
teori gangguan orde pertama, yang dipratayangkan di sini dan dibenarkan
dalam \cref{ch:b3:perturbation-theory}.
\end{proof}

\begin{omfigure}
\begin{tikzpicture}[font=\small, >=Stealth]
  \draw[->] (0,-3.0) -- (0,3.15);
  \node[anchor=south] at (0,3.15) {$z$};
  % sphere radius sqrt6
  \pgfmathsetmacro\R{2.449}
  \draw[gray, dashed] (0,0) circle (\R);
  \foreach \m/\c in {2/omDef, 1/omThm, 0/omExo, -1/omThm, -2/omDef} {
    \pgfmathsetmacro\zz{\m}
    \pgfmathsetmacro\rr{sqrt(\R*\R-\m*\m)}
    \draw[\c, thick, ->] (0,0) -- (\rr,\zz);
    \draw[gray!60, thin, dashed] (-0.05,\zz) -- (\rr,\zz);
    \node[anchor=west, \c] at ({\rr+0.1},\zz) {$m = \m$};
  }
  \node[anchor=east, font=\footnotesize] at (-0.15,2.0) {$L_z = m\hbar$};
  \node[anchor=north, font=\footnotesize] at (1.3,-2.75) {$\|\vect L\| = \sqrt{6}\,\hbar$};
\end{tikzpicture}

{\small ``Kuantisasi ruang'' untuk $l = 2$: vektor momentum sudutnya
berpanjang $\sqrt{l(l+1)}\,\hbar = \sqrt6\,\hbar$ tetapi hanya dapat
menawarkan kepada sumbu $z$ kelima proyeksi $m\hbar$ --- tak pernah
panjang penuhnya: karena komponennya tak serasi, $\vect L$ tak pernah
dapat terletak tepat sepanjang sumbu mana pun.}
\end{omfigure}

\begin{method}[Momentum sudut dalam praktik]\label{met:b3:quantum-angular-momentum:method}
(1) Labelilah keadaannya dengan $(l, m)$ --- atau $(j, m)$ untuk
aljabar abstraknya --- dan jangan pernah menanyakan dua komponennya
sekaligus. (2) Unsur matriksnya: pakailah rumus tangganya;
$\hat L_x = (\hat L_+ + \hat
L_-)/2$. (3) Potensial sentral: kutiplah $Y_l^m$, lalu selesaikanlah
hanya \omterm{prop:b3:schrodinger-three-dimensions:swave}{persamaan radialnya} beserta \omterm{thm:b3:quantum-angular-momentum:radial}{penghalang sentrifugalnya}. (4) Energi
rotasinya: $BJ(J+1)$, dengan garis pada $2B(J+1)$ --- tariklah $B$,
sehingga diperoleh panjang ikatannya, dari jaraknya yang sama. (5)
Pertanyaan magnetik: ubahlah momentum sudutnya menjadi momen sebesar
$\mu_{\text{B}}$ tiap $\hbar$, dan energinya sebesar $\mu_{\text{B}}B$
tiap satuan $m$.
\end{method}

\section{Latihan}

\begin{exercise}[$\star$]\label{exo:b3:quantum-angular-momentum:1}
(a) Turunkanlah $[\hat L_x, \hat L_y] = \iu\hbar\hat L_z$ dari
\omterm{def:b3:quantum-formalism:commutator}{komutator} kanoniknya. (b) Tunjukkanlah $[\hat L^2, \hat L_z] = 0$. (c)
Tunjukkanlah $[\hat L_z, \hat x] = \iu\hbar\hat y$: apa yang
dibangkitkan $\hat L_z$? (d) Mengapa ketiga komponennya tidak mempunyai
basis eigen bersama --- kecuali untuk satu keadaan tertentu (yang mana)?
\end{exercise}

\begin{exercise}[$\star$]\label{exo:b3:quantum-angular-momentum:2}
Untuk $l = 1$, dalam basis $\{\ket{1,1}, \ket{1,0}, \ket{1,-1}\}$:
(a) tuliskanlah matriks $\hat L_z$; (b) pakailah rumus tangganya untuk
menuliskan $\hat L_+$ dan $\hat L_-$; (c) susunlah $\hat L_x$ lalu
periksalah bahwa \omterm{def:b3:quantum-formalism:observable}{nilai eigennya} $\hbar, 0, -\hbar$; (d) keadaan dengan $L_x =
+\hbar$ diukur sepanjang $z$: berikanlah peluang ketiga hasilnya.
\end{exercise}

\begin{exercise}[$\star$]\label{exo:b3:quantum-angular-momentum:3}
Karbon monoksida: panjang ikatannya \qty{0.113}{nm}, massa tereduksinya
$\qty{1.14e-26}{kg}$. (a) Hitunglah $I$ dan $B = \hbar^2/2I$ dalam meV.
(b) Berapa frekuensi dan panjang gelombang garis $J = 1 \leftarrow 0$?
(c) Berapa dua garis berikutnya? (d) Seorang astronom mengukur jarak
sisirnya sampai lima angka: besaran molekul apa yang diperolehnya, dan
sampai ketelitian berapa?
\end{exercise}

\begin{exercise}[$\star$]\label{exo:b3:quantum-angular-momentum:4}
Sebuah partikel berada dalam keadaan $l = 2$. (a) Daftarkanlah hasil
yang mungkin dari pengukuran $L_z$. (b) Berapa sudut terkecil antara
$\vect L$ dan sumbu $z$, dengan memakai $\cos\theta = m/\sqrt{l(l+1)}$?
(c) Mengapa sudutnya tak pernah dapat nol? (d) Tunjukkanlah bahwa sudut
minimumnya menuju nol ketika $l \to \infty$: vektor klasiknya diperoleh
kembali.
\end{exercise}

\begin{exercise}[$\star\star$]\label{exo:b3:quantum-angular-momentum:5}
Bagian sudut $-\hbar^2\Delta$ adalah $\hat L^2$, dengan
\[
\hat L^2 = -\hbar^2\Big[\frac{1}{\sin\theta}\,\partial_\theta
(\sin\theta\,\partial_\theta) +
\frac{1}{\sin^2\theta}\,\partial_\varphi^2\Big] .
\]
(a) Periksalah bahwa $Y_1^0 \propto \cos\theta$ mempunyai \omterm{def:b3:quantum-formalism:observable}{nilai eigen}
$\hat L^2$ sebesar $2\hbar^2$. (b) Periksalah $Y_1^{\pm1} \propto \sin\theta\,
\eu^{\pm\iu\varphi}$ demikian pula, beserta \omterm{def:b3:quantum-formalism:observable}{nilai eigen} $\hat L_z$-nya.
(c) Periksalah paritasnya. (d) Mengapa $\braket{Y_1^0}{Y_1^1} = 0$ tanpa
menghitung integral apa pun?
\end{exercise}

\begin{exercise}[$\star\star$]\label{exo:b3:quantum-angular-momentum:6}
Garis rotasi manakah yang paling terang? Populasi aras $J$ adalah
$\propto (2J+1)\,\eu^{-BJ(J+1)/k_{\text{B}}T}$. (a) Jelaskanlah kedua
faktornya. (b) Tunjukkanlah bahwa maksimumnya terletak di dekat $J_{\max} \approx
\sqrt{k_{\text{B}}T/2B} - \tfrac12$. (c) Untuk CO pada \qty{10}{K}
(awan gelap) dan pada \qty{300}{K}: garis mana yang menguasainya? (d)
Sebaliknya: tangga CO teramati yang memuncak pada $J = 7$ membocorkan
suhu berapa?
\end{exercise}

\begin{exercise}[$\star\star$]\label{exo:b3:quantum-angular-momentum:7}
\omterm{prop:b3:quantum-angular-momentum:magnetic}{Efek Zeeman} normal. Sebuah garis spektrum pada \qty{500}{nm} berasal
dari transisi $l = 2 \to l = 1$ dalam medan $B = \qty{2}{T}$;
abaikanlah spinnya (berlaku untuk keadaan ``singlet'' khusus). (a)
Sketsalah aras yang terbelah itu. (b) Dengan kaidah pemilihan
$\Delta m = 0, \pm1$, tunjukkanlah bahwa hanya \emph{tiga} kedudukan
garis yang muncul, pada $0, \pm\mu_{\text{B}}B/h$.
(c) Hitunglah pembelahannya dalam GHz dan dalam pikometer panjang
gelombang. (d) Sebagian besar garis nyata terbelah menjadi lebih dari
tiga komponen (Zeeman ``anomali''): bahan yang hilang apa, yang dibawa
elektron itu sendiri, yang secara historis dibuktikannya?
\end{exercise}

\begin{exercise}[$\star\star$]\label{exo:b3:quantum-angular-momentum:8}
Dinding sentrifugal. Untuk hidrogen ($V = -e^2/4\pi\varepsilon_0r$),
tuliskanlah $U_{\text{eff}}$ untuk $l = 1$ dalam satuan $E_{\text{I}}$
dan $a_0$. (a) Tentukanlah minimumnya beserta nilainya. (b)
Tunjukkanlah $U_{\text{eff}} >
0$ untuk $r < \dots$ --- carilah jejari yang di dalamnya penghalangnya
menguasai. (c) Bandingkanlah $|\psi|^2$ di dekat $r = 0$ untuk keadaan
$s$ dan $p$ ($u \sim r^{l+1}$, diterima tanpa bukti): keadaan mana yang
``menyentuh'' intinya? (d) Penangkapan elektron --- sebuah inti yang
menelan salah satu elektron atomnya --- berlangsung nyaris seluruhnya
dari kulit $s$: jelaskanlah dalam satu kalimat.
\end{exercise}

\begin{exercise}[$\star\star$]\label{exo:b3:quantum-angular-momentum:9}
Pada $\ket{l, m}$: (a) tunjukkanlah $\langle\hat L_x\rangle = \langle\hat
L_y\rangle = 0$; (b) tunjukkanlah $\langle\hat L_x^2\rangle = \langle\hat
L_y^2\rangle = \tfrac12[l(l+1) - m^2]\hbar^2$; (c) periksalah \omterm{thm:b3:quantum-formalism:uncertainty}{hubungan
ketaktentuan} $\Delta L_x\,\Delta L_y \ge \tfrac\hbar2
|\langle\hat L_z\rangle|$ pada keadaan teregang $m = l$; (d)
tafsirkanlah gambaran ``kerucut'' pada gambarnya dalam terang (a) dan
(b).
\end{exercise}

\begin{exercise}[$\star\star\star$]\label{exo:b3:quantum-angular-momentum:10}
Pita getaran--rotasi. Sebuah transisi getaran inframerah
($\hbar\omega_0$) sebuah molekul dwiatom mengubah $J$ sebesar $\pm1$
secara serentak. (a) Tunjukkanlah bahwa garis serapannya jatuh pada $h\nu =
\hbar\omega_0 + 2B(J+1)$ (cabang $R$, dari $J \to J+1$) dan
$h\nu = \hbar\omega_0 - 2BJ$ (cabang $P$, $J \to J-1$), $J \ge
1$. (b) Mengapa ada \emph{celah} tepat pada $\hbar\omega_0$? (c)
Sketsalah pitanya: dua sisir yang mengapit satu garis pusat yang hilang,
dengan kekuatan tiap garisnya mengikuti
\cref{exo:b3:quantum-angular-momentum:6}. (d) Pita
\qty{15}{\micro m} pada \cref{pb:b3:harmonic-oscillator:1} mempunyai
anatomi persis seperti ini: apa yang menetapkan lebar keseluruhannya,
dan karenanya ``sayap'' yang membuat karbon dioksida tambahan berdaya guna?
\end{exercise}

\begin{exercise}[$\star\star\star$]\label{exo:b3:quantum-angular-momentum:11}
Menuntaskan aljabarnya. (a) Dari $\hat J_\mp\hat J_\pm = \hat J^2 -
\hat J_z^2 \mp \hbar\hat J_z$, hitunglah $\|\hat J_\pm\ket{j,m}\|^2$
beserta koefisien tangganya. (b) Tunjukkanlah bahwa jika tangganya gagal
berhenti, suatu norma akan menjadi negatif: tunjukkanlah keadaan yang
melanggarnya.
(c) Simpulkanlah bahwa $m_{\max} - m_{\min} = 2j$ haruslah bilangan
bulat tak negatif lalu daftarkanlah $j$ yang diizinkan. (d) Di manakah
tepatnya argumen itu \emph{gagal} menyingkirkan nilai setengah bulat ---
dan syarat tambahan yang mana (ketunggalan nilai pada bolanya) yang
menyingkirkannya hanya bagi gerak orbital?
\end{exercise}

\begin{exercise}[$\star\star\star$]\label{exo:b3:quantum-angular-momentum:12}
Einstein--de Haas: momentum sudut itu mekanis. Sebuah silinder besi
(jejari $a = \qty{1.0}{cm}$, massa jenis $\rho = \qty{7.9e3}{kg/m^3}$,
$n = \qty{8.5e28}{atom/m^3}$) tergantung pada serat puntir; membalikkan
kemagnetannya membalikkan sekitar satu $\hbar$ momentum sudut tiap
atomnya. (a) Menurut kekekalannya, batangnya harus terhentak: hitunglah
momentum sudut tiap satuan volume yang dipindahkan. (b) Dengan momen kelembaman batangnya $\tfrac12 M
a^2$, tunjukkanlah bahwa batangnya memperoleh $\omega = 2n\hbar/\rho a^2$
lalu hitunglah. (c) Percobaan 1915 menggerakkan pembalikannya pada
resonansi seratnya agar hentakan mungilnya menumpuk: perkirakanlah
amplitudo osilasinya setelah $Q \sim 100$ pembalikan resonan, dengan
memakai $\omega$ tiap hentakan yang Anda peroleh. (d) Nisbah terukur
momen terhadap momentum sudutnya ternyata dua kali ramalan orbitalnya
$e/2m_{\text{e}}$: faktor dua itu mengumumkan apa?
\end{exercise}

\section{Soal: Menimbang galaksi pada 2,6 milimeter}

\begin{problem}[{Soal akhir pekan --- karbon monoksida, astronomi
radio dan semesta yang dingin}]\label{pb:b3:quantum-angular-momentum:1}
Sebagian besar materi pembuat bintang sebuah galaksi adalah hidrogen
molekul yang dingin --- dan ia tak terlihat: H$_2$ yang setangkup tidak
mempunyai dipol dan tidak mempunyai garis rotasi pada suhu awan yang
dingin. Penyelesaian astronomi adalah pendampingnya yang setia. Satu
molekul CO tiap $\num{e4}$ H$_2$, dengan tangga \qty{115}{GHz}-nya,
menyalakan setiap awan molekul di langit. Data: untuk CO, $B/h = \qty{57.6}{GHz}$ ($B =
\qty{0.238}{meV}$); panjang ikatannya $r_0 = \qty{0.113}{nm}$; massa
tereduksinya $\mu = \qty{1.14e-26}{kg}$; $k_{\text{B}}T$ pada \qty{10}{K}
adalah \qty{0.862}{meV}.

\medskip\noindent\textbf{Bagian I --- Rotor kuantumnya.}
\begin{enumerate}
  \item Dari $\hat H = \hat L^2/2I$, berikanlah energinya $E_J$
        beserta degenerasinya.
  \item Periksalah $B = \hbar^2/2I$ untuk CO dari $r_0$ dan $\mu$.
  \item Dengan kaidah pemilihan $\Delta J = \pm1$, turunkanlah
        frekuensi serapannya $\nu_{J+1\leftarrow J} = 2B(J+1)/h$
        lalu daftarkanlah tiga yang pertama.
  \item Mengapa garisnya \emph{berjarak sama} --- dan apa yang
        diserahkan jarak terukurnya kepada astronom (lewat $I$)?
  \item Hitunglah panjang gelombang garis dasarnya, lalu jelaskanlah
        mengapa mengamatinya memerlukan tempat yang tinggi dan kering
        (apa yang menyerap gelombang milimeter di atmosfer kita
        sendiri --- ingatlah tetangga molekul pada
        \cref{pb:b3:harmonic-oscillator:1}, yaitu air)?
  \item Foton garis $1 \to 0$ membawa $\hbar$ momentum sudut: mengapa
        $\Delta J = \pm1$ bukan sekadar patokan kasar melainkan sebuah
        hukum kekekalan?
\end{enumerate}

\medskip\noindent\textbf{Bagian II --- Mengapa CO dan bukan H$_2$.}
\begin{enumerate}[resume]
  \item H$_2$ beratom sejenis: nyatakanlah mengapa ia sama sekali
        tidak memancarkan radiasi dipol rotasi.
  \item H$_2$ juga \emph{ringan}: hitunglah tetapan rotasi H$_2$
        ($\mu = m_{\text{p}}/2$, $r_0 = \qty{0.074}{nm}$)
        beserta energi transisi terjangkau pertamanya; bandingkanlah
        dengan $k_{\text{B}}T$ pada \qty{10}{K}: bahkan lewat pintu
        belakang kuadrupol, dapatkah H$_2$ yang dingin memancar?
  \item Anak tangga pertama CO berharga $2B = \qty{0.48}{meV}$
        berbanding $k_{\text{B}}T = \qty{0.86}{meV}$: apakah tangganya
        hidup pada \qty{10}{K}?
  \item Hitunglah nisbah populasinya $n_1/n_0$ pada \qty{10}{K}
        (termasuk degenerasinya).
  \item Di dekat $J$ berapa populasinya memuncak pada \qty{10}{K}
        (pakailah $J_{\max} \approx \sqrt{k_{\text{B}}T/2B} - \tfrac12$)?
  \item Ringkaslah dalam satu kalimat perjanjian dagangnya: apa yang
        disediakan CO, dan apa yang harus dianggap tentang nisbah CO
        terhadap H$_2$.
\end{enumerate}

\medskip\noindent\textbf{Bagian III --- Membaca langitnya.}
\begin{enumerate}[resume]
  \item Garis sebuah awan tiba pada \qty{115.156}{GHz} alih-alih
        \qty{115.271}{GHz} di laboratorium: hitunglah kecepatan awannya
        sepanjang garis pandang, beserta arahnya.
  \item Garisnya \emph{melebar} secara Doppler oleh gerak internalnya
        sebesar $\pm\qty{2}{km/s}$: hitunglah lebar garisnya dalam MHz,
        lalu bandingkanlah dengan lebar termal yang diharapkan pada
        \qty{10}{K} ($v_{\text{th}} \sim \sqrt{k_{\text{B}}T/m_{\text{CO}}}
        \approx \qty{55}{m/s}$): apa yang menguasai pelebarannya?
  \item Memetakan pergeseran Doppler CO melintasi cakram sebuah galaksi
        spiral menghasilkan kurva rotasinya $v(r)$. Kurvanya tetap
        \emph{datar} jauh melampaui cakram tampaknya: ingatlah (jilid
        Tahun 1, gravitasi) apa yang seharusnya dilakukan $v(r)$ di
        sekitar massa yang terpusat, lalu nyatakanlah apa yang
        disiratkan kedatarannya.
  \item Nisbah keterangan garis $2\to1$ terhadap garis $1\to0$
        mengukur populasi arasnya: besaran fisis tunggal apa dari awan
        itu yang diserahkannya?
  \item ALMA memilah CO dalam galaksi yang cahayanya berangkat ketika
        semesta masih muda; frekuensi teramati garis $1 \to 0$ dari
        sebuah galaksi pada ``pergeseran merah $z = 1$'' tiba pada
        separuh nilai diamnya --- pada frekuensi berapa, dan
        (dengan mengingat \cref{ex:b3:relativistic-kinematics:redshift})
        apa yang meregangkannya?
  \item Dari luminositas garis CO yang terukur, para astronom
        menyebutkan massa awan dalam satuan massa matahari:
        daftarkanlah rantai konversinya (foton garis $\to$ cacah CO
        $\to$ massa H$_2$) beserta mata rantainya yang terlemah.
\end{enumerate}

\medskip\noindent\textbf{Bagian IV --- Ruang bayi yang dingin.}
\begin{enumerate}[resume]
  \item Teras pembentuk bintang berada pada \qty{10}{K}: menurut hukum
        Wien pijar termalnya memuncak di dekat \qty{290}{\micro m} ---
        tak terlihat oleh teleskop optis mana pun. Dalam satu kalimat:
        mengapa tangga rotasinya merupakan \emph{satu-satunya} termometer
        dan penggaris yang ditawarkan teras semacam itu?
  \item Aras $J = 1$ terletak \qty{5.5}{K} (dalam satuan suhu) di atas
        keadaan dasarnya: mengapa hal itu menjadikan garis $1\to0$
        termometer yang sangat halus justru pada rezim \qty{10}{K}
        (alih-alih, misalnya, sebuah garis optis pada
        \qty{2}{eV} $\sim \qty{23000}{K}$)?
  \item Mengapa teras yang \emph{meruntuh} akhirnya tidak lagi terlihat
        dalam CO ($1\to0$) --- tinjaulah apa yang dilakukan kerapatan
        tinggi dan debu terhadap garisnya dan terhadap lolosnya.
  \item Awan molekul juga menunjukkan garis NH$_3$, HCN, dan berpuluh
        lainnya, masing-masing dengan $B$ dan dipolnya sendiri: apa yang
        dapat diurai para astronom berkat kekayaan ``pemutar gelombang
        radio kimia'' ini?
  \item Seluruh langit berpijar pada \qty{2.7}{K} (latar gelombang
        mikro kosmik): apa yang terjadi pada ketegasan garis CO sebuah
        awan ketika suhu awannya sendiri mendekati
        \qty{2.7}{K}, dan mengapa hal itu merupakan lantai yang keras
        bagi termometer ini?
  \item Cakram milimeter harus mempertahankan bentuknya sampai sekitar
        $\lambda/20$: hitunglah kelonggaran itu pada \qty{115}{GHz},
        lalu bandingkanlah dengan kelonggaran yang harus dipenuhi cermin
        optis (\qty{500}{nm}) --- mengapa cakram dua belas meter milik
        ALMA dapat dibangun di ruang terbuka?
  \item Ringkaslah hasil bernama itu: sebuah tangga $BJ(J+1)$ dengan $B/h =
        \qty{57.6}{GHz}$, yang terisi pada sepuluh kelvin dan dibaca pada
        \qty{2.6}{mm}, adalah cara massa, suhu, kecepatan dan rotasi
        semesta yang dingin --- beserta bukti adanya materi gelap di
        sekitar galaksi spiral --- diukur.
\end{enumerate}
\end{problem}
